\documentclass[11pt]{article}
\usepackage{amsmath,amssymb,amsthm,hyperref}
\newtheorem{theorem}{Theorem}\newtheorem{proposition}[theorem]{Proposition}
\newcommand{\D}{\mathsf D}\newcommand{\E}{\mathcal E}\newcommand{\C}{\mathcal C}\newcommand{\F}{\operatorname{Fam}}
\title{A directed bridge from conditional native bind to all-coupling span feedback}
\date{16 November 2026}
\begin{document}\maketitle
\paragraph{Status.} The written bridge, global writer lift and feedback-descent obstruction have passed independent review at the scopes stated below. This is a relative construction, not a claimed conservative strong completion of the original context/content crossing. It retains actual coefficient arrows and derives a coherent directed comparison from conditional bind to joint coupling. It changes the interconnection interface explicitly: a span retains all compatible joint laws; conditional bind selects a particular law. No novelty, formal certification, arbitrary polynomial content, or infinite semantics is claimed.

\section{Source and target}
Let $\mathcal V$ be a locally small symmetric monoidal category with the finite limits needed below. Put $\E=\mathcal V/I$, with its induced augmented tensor $\otimes$ and unit $(I,\mathrm{id})$. This unit is terminal in $\E$. Thus there are canonical natural arrows $C\otimes D\to C,D$, formed using their augmentations. The induced tensor need not be the categorical product. Alternatively any finitely complete symmetric monoidal category with terminal tensor unit can be used as $\E$.

Let $\D X$ be normalized nonnegative probability laws on a finite raw set $X$, with zero-weight points retained. Define $\C$ to have objects $(W,p,(C_w))$, where $W$ is finite, $p\in\D W$, and $C_w\in\E$. A native arrow is an actual function $r:W\to W'$ and actual coefficient arrows $C_w\to C'_{r(w)}$ satisfying $r_*p=p'$. The previously established native multilimit construction gives finite limits in $\F(\C)$, where all small families, not merely finite families, are admitted. Let $\mathsf S=\operatorname{Span}(\F(\C))$ be the bicategory with maps of spans as 2-cells. Its tensor is categorical product in $\F(\C)$. For an object $X$, the cup is the span $1\leftarrow X\xrightarrow{\Delta}X\times X$ and the cap is $X\times X\xleftarrow{\Delta}X\to1$. Pullback uniqueness gives their snake isomorphisms. This is the standard compact span construction; see Stay, \emph{Compact closed bicategories}, Theory and Applications of Categories 31 (2016), 755--798, \url{https://www.tac.mta.ca/tac/volumes/31/26/31-26.pdf}. The locally discrete specialization retains actual maps of spans.

The source bicategory $\mathsf K$ has finite sets as objects. An object $P:A\to B$ of a hom category consists, for every $a$, of a finite raw witness set $W_a$, law $\mu_a$, coefficients $C_{a,w}\in\E$, and output $b(a,w)\in B$. A 2-cell consists of actual maps $r_a$ preserving pushforward laws and outputs, with actual coefficient arrows. Composition $Q P$ has witnesses $(w,z)$, weight $\mu_a(w)\nu_{b(a,w)}(z)$ and coefficient $C_{a,w}\otimes D_{b(a,w),z}$. Identities have singleton witness and tensor unit. This is the finite normalized augmented native-context fragment. This initial theorem omits global writer content. The subsequent global-writer lift retains it in full; arbitrary original polynomial contentads still require a separate argument.

For a finite set $A$, define
\[
 L A=((A,p,(I)_a))_{p\in\D A}.
\]
For a process $P$, define a family $\Gamma P$, indexed by $p\in\D A$, whose component has raw set $U_P=\coprod_a W_a$, law $p(a)\mu_a(w)$, and coefficient $C_{a,w}$. Its two native legs give a span
\[
 LP=(LA\xleftarrow{s}\Gamma P\xrightarrow{t}LB).
\]
The source index map is $p\mapsto p$; the target index map is $p\mapsto K_Pp$, the pushforward along $b$. Both leg coefficient arrows are augmentations. Native process cells induce span maps using the same witness and coefficient arrows for every prior $p$.

\section{The directed composition theorem}
\begin{theorem}
The preceding construction gives a locally faithful oplax functor with invertible unit constraints
\[
 L:\mathsf K\longrightarrow\mathsf S,
 \qquad \chi_{Q,P}:L(QP)\Longrightarrow LQ\circ LP.
\]
Here ``oplax'' fixes the displayed direction, regardless of naming conventions. The object assignment has coherent isomorphisms $L(A\times B)\cong LA\times LB$ and $L1\cong1$. There are compatible oplax tensor comparisons $L(P\otimes Q)\to LP\otimes LQ$. Neither the composition nor tensor comparisons are generally invertible. Local fullness is not asserted.
\end{theorem}
\begin{proof}
The pullback apex of $LP,LQ$ is the native multilimit family. Its indices consist of a prior $p$ on $A$, the forced intermediate prior $q=K_Pp$ on $B$, and a coupling on
\[
 U_P\times_B U_Q=\{((a,w),(b,z)):b=b(a,w)\}
\]
with the required marginals. Its coefficient at a tuple is $C_{a,w}\times D_{b,z}$ in $\E$, since the shared boundary coefficient is terminal. The law
\[
 \lambda_p(a,w,z)=p(a)\mu_a(w)\nu_{b(a,w)}(z)
\]
has precisely those two marginals. Therefore it selects a component of this family. The canonical raw bijection from dependent-pair witnesses and the canonical coefficient arrow $C\otimes D\to C\times D$ give $\chi_{Q,P}$. This is an actual native arrow, not only a law comparison.

Naturality follows on raw witnesses from the coordinate maps, on laws by summing the product weights over their fibres, and on coefficients from naturality of the augmentation-derived projections. For three processes, both composite comparison routes select the same complete law
\[
 p(a)\mu_a(w)\nu_{b(a,w)}(z)\xi_{c(b,z)}(v)
\]
on the same retained triple of witnesses. After the canonical pullback reassociation, both coefficient arrows have the same three projections, namely tensoring with the augmentations of the other two factors. Symmetric monoidal coherence and the augmentation equations make these projections identical. Product uniqueness then proves equality of the actual coefficient arrows. The family multilimit characterization fixes the family index as well. This proves the oplax associativity equation. Singleton identity witnesses and terminal-unit projections prove both unit equations; $L$ takes identities to identity spans up to the displayed native isomorphism.

For local faithfulness, when $A$ is nonempty choose any prior $p$. Its raw apex includes every $a,w$, even when $p(a)=0$. Equality of induced span maps therefore implies equality of every original witness and coefficient map. If $A$ is empty the source hom has only the unique empty data and the assertion is immediate.

The product $LA\times LB$ has one component for each triple $(p,q,\lambda)$ where $\lambda$ couples $p,q$ on $A\times B$. Marginals uniquely determine $p,q$ from $\lambda$, so reindexing by $\lambda$ and using the terminal-coefficient isomorphism gives $L(A\times B)\cong LA\times LB$. Its finite-product coherence follows by identifying each component with its complete joint law and each raw witness with its tuple.

For the tensor comparison, a prior $r$ on $A\times A'$ gives the joint witness law
\[
 r(a,a')\mu_a(w)\nu_{a'}(z).
\]
Its marginals are the two apex laws indexed by the marginals of $r$. Again use $C\otimes D\to C\times D$. The symmetry, tensor associativity and composition--tensor compatibility equations follow by flattening to the same addressed tuple, the same product of conditional weights, and the same coefficient projections. In particular this argument does not assert that arbitrary independent resampling equals a shared coupling.
\end{proof}

The comparison is usually noninvertible. Let $P$ generate a fair bit $w$ but return only a constant intermediate output, and let $Q$ generate a fair bit $z$. The bind-selected joint law is independent. The span composite also has diagonal and antidiagonal fair laws as distinct components. No inverse can make the family indices agree. Even if the laws were unique, $C\otimes D\to C\times D$ could be noninvertible.

Local fullness fails independently. With two inputs, a constant returned output, terminal context coefficients, and an internal fair-bit witness at each input, a span endomorphism can swap that bit for some input priors and leave it fixed for others. It preserves both span legs and each prior's law. It need not arise from one prior-independent family $r_a$. Thus the construction preserves distinctions without claiming that its target adds no comparisons between old images.

\section{Exact feedback meaning}
Take $P:A\times S\to B\times S$ and use the object-product isomorphisms above. The compact span trace has apex the equalizer of the two state legs of $\Gamma P$. Its components are exactly the priors $p\in\D(A\times S)$ such that
\[
 p(a,s)\mu_{a,s}(w)>0\ \Longrightarrow\ s_{\rm out}(a,s,w)=s.
 \tag{*}
\]
The raw carrier in such a component is the full equality subset, including its zero-weight witnesses, with restricted probability law and unchanged coefficients. Nonclosed zero-weight raw points are not retained by this equalizer. The outside legs retain input $a$ and output $b$.

To verify the assertion, a cone into the equalizer must push forward to the entire original apex law. The underlying raw equalizer inclusion forces zero mass outside its image; conversely this condition gives its unique normalized restricted law. The state-index equality follows from this law, so imposes no further condition. Coefficient equalizers of the two augmentation arrows to the terminal coefficient are identities. The native family multilimit construction then supplies exactly the stated apex.

This is pathwise fixed-loop feasibility, not merely stationarity. A fair-bit toggle has a stationary uniform law but has no equality witness and hence an empty feedback apex. Nor is this conditioning: mass outside the equality event is not discarded and renormalized. The operation is total as a span construction, allowing an empty apex, while existence of a nonempty component is the derived admissibility condition (*). No selected prior or selected coupling is part of the feedback operation.

\section{A strict interconnection separation}
Let $S=\{0,1\}$. Both $P:S\to S$ and $Q:S\to S$ discard their input and emit a fresh fair bit. Their ordinary conditional bind also emits a fair bit independently of its initial state. For every prior $p$ on initial $s$, the closed-event mass is $1/2$, so the equalizer trace of $L(QP)$ is empty.

In contrast, $LQ\circ LP$ admits the joint law
\[
 \Pr(s,t,u)=\begin{cases}1/4&u=s,\\0&u\ne s.\end{cases}
\]
Its $(s,t)$ marginal is the first channel's joint law for a uniform input, and its $(t,u)$ marginal is the second channel's joint law for a uniform intermediate state. Hence it is an admitted all-coupling component. Every positive-weight witness has $u=s$, so its feedback equalizer is nonempty. This is a concrete new admissible interconnection produced by retaining shared innovations across the cut. It is not a claim that fresh independent resampling itself closes; the directed comparison marks exactly the independent component and does not become an equivalence after closure.

The comparison $L(QP)\to LQ\circ LP$ transports through span trace because trace acts functorially on span maps. Thus the empty/nonempty separation is compatible with, rather than a counterexample to, the directed bridge. Finite branching is represented by arbitrary finite product interfaces and the all-coupling product of spans; its tensor comparison retains conditional independence as one specific admitted component. Arbitrary original polynomial substitution is not inferred from this representation.

\section{What this does and does not settle}
This supplies one coherent relative family: normalized conditional native bind maps directionally into all-coupling span composition; joint assembly and fixed-loop feedback are generated by the same finite-limit calculus. Its feedback does not require a weak product in the original exact-law category because the target retains a family of alternative laws. The result neither constructs a selected-law retraction nor contradicts the selected-law-shadow obstruction.

The governing v0.2 target remains open. The subsequent lift supplies global writer grades without discarding unused labels. Required further work includes relating this interface to the independent context/content comma presentation and arbitrary polynomial substitution; determining an appropriate universal property for the process bridge rather than just citing the ordinary span construction; and separating infinite-run semantics from these finite fixed-loop witnesses. Retained-state execution and PR77's unnormalized instantaneous trace are not identified with this operation. A strong conservative embedding preserving old bind is specifically not established: the noninvertible compositor is a declared and mathematically necessary change of preservation judgment in this construction.
\appendix
\section{Independent global writer data}
Let $\mathcal D$ be any locally small symmetric monoidal category, independently of the augmented context-coefficient category $\mathcal E$. Write $\mathsf W=\operatorname{Fam}(\mathcal D)$ for the category of small families. A global writer object is
\[
 Q=(O,(D_o)_{o\in O}).
\]
A writer morphism is one function $v:O\to O'$ and one actual arrow $D_o\to D'_{v(o)}$ for every $o$. The writer tensor $Q\boxtimes R$ has label set $O\times O_R$ and coefficient $D_o\otimes D^R_r$. Its unit has one label and the tensor unit of $\mathcal D$. There is no condition that all labels occur in the body, no deletion of unused labels, and no identification of arrows invisible to the label function. Small infinite label families are allowed even though each process row has finite witnesses.

A source process now consists of the previous native context data, one global $Q$, and body
\[
 (a,w)\longmapsto (b(a,w),o(a,w))\in B\times O.
\]
A source comparison carries the previous witness and context-coefficient arrows and one full writer morphism $v:Q\to Q'$ satisfying the body square. This same $v$ acts at every input and witness. Composition retains exactly $Q\boxtimes R$, with the paired label $(o(a,w),r(b(a,w),z))$. It does not replace the global writer with the family of realized labels. Thus unused writer coefficients survive composition and feedback. Independent tensor has the same product-label writer and uses the symmetry of $\mathcal D$ in interchange.

\section{Writer-decorated spans}
Put $\mathcal B=\operatorname{Fam}(\mathcal C)$, where $\mathcal C$ is the native finite-probability context category of the preceding construction. There is a total raw-carrier functor
\[
 U:\mathcal B\to\mathbf{Set},\qquad U((X_i)_{i\in I})=\coprod_{i\in I}|X_i|.
\]
No preservation of finite limits by $U$ is asserted. In fact its product can carry several coupling indices over the same pair of raw points.

A decorated span from $A$ to $B$ is a span $A\leftarrow X\to B$ in $\mathcal B$, one global writer $Q$, and a function $\ell:UX\to O$. A 2-cell is a map of spans $h:X\to X'$ and one full writer arrow $v:Q\to Q'$ such that
\[
 \ell'\,Uh=|v|\,\ell.
\]
Composition of the underlying spans uses a chosen pullback $X\times_B Y$. The global writer is $Q\boxtimes R$ and its label is
\[
 z\longmapsto(\ell_X(U\pi_X z),\ell_Y(U\pi_Y z)).
\]
The identity span has unit writer and constant label. Tensor of spans uses products in $\mathcal B$, tensor of writers uses $\boxtimes$, and the label is obtained from the two product projections. Vertical composition of cells composes the actual global writer arrows; horizontal composition uses their writer tensor and the induced pullback map.

\begin{theorem}
These data form a symmetric monoidal bicategory of writer-decorated spans. Its objects are canonically self-dual, using the ordinary diagonal/terminal span cups and caps decorated by the unit writer. Its trace restricts the old span apex to the state equalizer, retains the entire global writer unchanged, and restricts only its label function. Trace sends a 2-cell $(h,v)$ to the induced equalizer map together with the same full writer arrow $v$.
\end{theorem}
\begin{proof}
Pullback reassociation has the same projections to every original apex. Consequently its induced labels are the same ordered label tuple, up to the corresponding writer associator. This proves that it is a decorated-span isomorphism. Unit pullback isomorphisms use writer unitors. Pentagon and triangle commute on span apices by pullback uniqueness and on writers by monoidal coherence; projection formulas prove the label equation. The same argument handles tensor associativity and symmetry. The interchange isomorphism reorders the middle two writer factors using their symmetry and reorders the pullback tuple by its canonical isomorphism. Thus it is valid on actual writer coefficient arrows, not merely labels. Naturality follows from the pullback projection equations and naturality in $\mathsf W$.

Decorate both cup and cap by the writer unit. Each snake composite has writer a tensor of units and canonical unit label; the usual span snake isomorphism and the writer unit coherence give a decorated-span isomorphism. Their coherence follows from the same two independent uniqueness/coherence arguments. The composite defining trace contains the original writer and only units. Applying its canonical unitors leaves precisely $Q$, including every unused coefficient. The ordinary span trace apex is the equalizer of the two state legs; its label is the old label precomposed with the equalizer inclusion. Its underlying native equalizer is computed by the general family multilimit construction; for process-image spans over $LS$, the explicit prior/support criterion above applies. For arbitrary state objects, the coefficient equalizers can be nontrivial and must not be replaced by identities. A map of decorated spans preserves the two state legs and its writer-label equation restricts to this equalizer, proving the claimed action on cells. None of these arguments replaces $U$ of a limit by a limit of $U$.
\end{proof}

\section{Lift of the oplax bridge}
The image of a source process retains exactly its global $Q$. At the component of $\Gamma P$ indexed by input prior $p$, put
\[
 \ell_P(p,a,w)=o(a,w).
\]
A native source cell induces the old apex map and retains its single full writer map. The composition comparison from the preceding construction retains the global writer $Q\boxtimes R$ identically; on its selected Markov component the label of the dependent-pair witness is exactly the paired label prescribed by the pullback projections. The tensor comparison likewise has identity global writer and the same paired label. Units and all coherence cells carry the appropriate canonical writer unitors, associators and symmetries.

\begin{theorem}
The locally faithful oplax bridge with invertible unit constraints and its compatible oplax tensor comparisons lift to the full global-writer source and writer-decorated span target. For nonempty input interfaces, local faithfulness follows from the old raw/context argument and retention of the full writer map. For empty input interfaces it still holds: the apex and label data are empty, but the full writer object and every writer arrow remain part of the target hom category.
\end{theorem}
\begin{proof}
The displayed label equation proves well-definedness of the lifted compositor. In all composition and tensor coherence diagrams, the writer maps are precisely those of the same monoidal expressions in the original $Q,R,\ldots$. Their equality is symmetric monoidal coherence in $\mathsf W$; the apex equations are the already proved equations from the preceding construction. Hence the composite decorated equations commute. An equality of target cells includes equality of the writer maps as arrows of $\mathsf W$, even if the apex is empty or some labels are unused. This supplies the missing writer part of faithfulness in every case. The strict interconnection separation of the preceding construction uses the unit writer and therefore survives this lift. More generally it survives any chosen nonempty global writer with admissible body labellings and arbitrary unused labels and coefficients: emptiness/nonemptiness of the effect apex is unchanged and no label is erased.
\end{proof}

\section{Pure-endpoint squares}
This construction also respects changing pure endpoints, not just fixed hom categories. A pure map $f:A\to A'$ acts on $LA$ by the family map with index $p\mapsto f_*p$, raw map $f$, and identity terminal-coefficient maps. For pure $f:A\to A'$ and $g:B\to B'$, a source square from $P$ to $P'$ consists of maps
\[
 r_a:W_a\to W'_{f(a)},\qquad (r_a)_*\mu_a=\mu'_{f(a)},
\]
actual context-coefficient maps, one global writer arrow $v$, and the output/label square. Its induced apex map sends
\[
 (p,a,w)\longmapsto(f_*p,f(a),r_a(w)).
\]
At each target raw point $(a',w')$, its pushed-forward mass is
\[
 \sum_{a:f(a)=a'}p(a)\sum_{w:r_a(w)=w'}\mu_a(w)
 =\sum_{a:f(a)=a'}p(a)\mu'_{a'}(w')
 =(f_*p)(a')\mu'_{a'}(w').
\]
Thus it is a genuine native map even if $f$ identifies inputs or is not surjective. Context coefficient maps are the supplied ones and the global writer map remains $v$. Its two leg equations are the pure source/output equations. Identity and vertical composition hold on these actual maps. Horizontal comparison naturality is the same dependent-product pushforward calculation as in the preceding theorem, with the intermediate endpoint square ensuring that continuation rows are compared at their correct target input. The writer square follows from its one globally supplied map.

Accordingly there is an oplax horizontal double-functor formulation with pure finite-set maps vertically and maps of decorated spans as squares. This formulation does not introduce arbitrary independent prior-indexed source maps. The span target still has additional cells, so local fullness remains false. Its joint composition and feedback remain the declared alternative-coupling interface, not a strong preservation of the original conditional bind.

\section{A descent obstruction between the two feedback interfaces}
\label{sec:dilation-descent}

The atomic-generation theorem supplies normalized finite-wire presentations of bounded native one-cells. Such presentations cannot be discarded when transporting the all-coupling span interpretation through feedback. The obstruction already occurs for finite raw sets, terminal context and writer coefficient categories, and the unit global writer family. In particular it is not caused by nontrivial coefficient arrows, unused-label deletion, an infinite carrier, or a change of endpoint type.

Write $\operatorname{Tr}^{\mathrm{fin}}$ for the native finite-measure feedback that retains all loop-consistent raw witnesses, without normalization. Write $L$ for the normalized process-to-span bridge and $\operatorname{Tr}^{\mathrm{span}}$ for compact span feedback. For a normalized $F:U\to U$, a component of the latter is a prior $p$ on $U$ whose complete joint law $p(u)\mu_u(w)$ is supported on $u_{\rm out}(w)=u$. This is the previously derived whole-law admissibility criterion, not postselection.

\paragraph{Proposition (failure of normalized-dilation descent).}
There are normalized processes $F_{\rm pack},F_{\rm split}:U\to U$ on the same two-element loop set, with the same unit global writer family, for which
\[
 \operatorname{Tr}^{\mathrm{fin}}_U(F_{\rm pack})
 \cong \operatorname{Tr}^{\mathrm{fin}}_U(F_{\rm split})
\]
as full native one-cells, but
\[
 \operatorname{Tr}^{\mathrm{span}}_{LU}(LF_{\rm pack})
 \not\cong
 \operatorname{Tr}^{\mathrm{span}}_{LU}(LF_{\rm split}).
\]
The first span apex is nonempty and the second is the empty family. Consequently no function on native isomorphism classes of bounded finite-measure one-cells can assign the span-trace isomorphism class by choosing an arbitrary normalized dilation and applying $L$ followed by span feedback.

\begin{proof}
Take external input and output $1$ and loop set $U=\{0,1\}$. All context coefficients are the unit, and the full global writer family is the singleton family $Q=(\{*\},I)$. Every body returns writer label $*$; this same whole grade is retained throughout. The process $X:1\to1$ to be presented has precisely two raw witnesses $x_0,x_1$, each of weight $1/2$.

The packed dilation has the following raw rows. Every displayed witness is retained, including the zero-weight dummy:
\[
\begin{array}{c|c|c|c}
\text{input loop}&\text{raw witness}&\text{weight}&\text{output loop}\\\hline
0&x_0&1/2&0\\
0&x_1&1/2&0\\
0&d_0&0&1\\
1&d_1&1&0
\end{array}
\qquad (F_{\rm pack}).
\]
Both rows have total mass one. Its only loop-consistent raw witnesses are $(0,x_0)$ and $(0,x_1)$. In particular the zero-weight dummy $d_0$ is off-loop, so native feedback removes it by the equality condition, not by deleting zero-weight witnesses as a convention. The other dummy is also off-loop. The raw bijection $(0,x_i)\mapsto x_i$, equipped with identity context-coefficient maps and the identity on the entire $Q$, identifies its finite-measure trace with $X$.

The split dilation has rows
\[
\begin{array}{c|c|c|c}
\text{input loop}&\text{raw witness}&\text{weight}&\text{output loop}\\\hline
0&x_0&1/2&0\\
0&e_0&1/2&1\\
1&x_1&1/2&1\\
1&e_1&1/2&0
\end{array}
\qquad (F_{\rm split}).
\]
Again both rows are normalized. Its loop-consistent witnesses are exactly $(0,x_0)$ and $(1,x_1)$. The raw bijection $(i,x_i)\mapsto x_i$, with the same identity coefficient and full writer maps, identifies this finite-measure trace with the same $X$. These are native isomorphisms, not merely equal output distributions.

For $F_{\rm pack}$, the prior $p=\delta_0$ puts the entire joint probability law on the two closed witnesses. It therefore supplies a span-feedback component whose raw equality carrier is those two witnesses with weights $1/2,1/2$. More precisely, the mass outside equality is $p(1)$, so $\delta_0$ is its only admitted prior.

For $F_{\rm split}$, every input row puts mass $1/2$ outside equality. For every normalized prior $p$,
\[
 \Pr_p[u_{\rm out}\ne u]
 =\tfrac12p(0)+\tfrac12p(1)=\tfrac12.
\]
No prior satisfies whole-law loop admissibility, so the span-feedback apex is the empty family. The uniform prior is stationary for the output transition, but does not satisfy the stronger pathwise equality condition. No isomorphism in the family category can identify an empty index family with a nonempty one; consequently these scalar spans cannot be isomorphic. The writer remains $Q$ even on the empty apex, so the separation does not rely on forgetting a global grade in an empty case.

If an assignment on native isomorphism classes descended the operation $F\mapsto\operatorname{Tr}^{\mathrm{span}}(LF)$ through $F\mapsto\operatorname{Tr}^{\mathrm{fin}}(F)$ for all normalized finite-wire presentations, its value on $[X]$ would have to be both displayed nonisomorphic span classes. This is impossible.
\end{proof}

\paragraph{Corollary (no invertibly feedback-preserving extension of the normalized bridge).}
There is no extension of the normalized bridge $L$ to the finite-measure native doctrine which preserves native isomorphisms, agrees with $L$ on normalized processes, and preserves finite-wire feedback by invertible comparisons, with the stated object/product identifications.
\begin{proof}
The common process $X$ above is itself normalized: its two weights sum to one. Therefore agreement on normalized processes fixes its image to $LX$, up to the agreed isomorphism. This scalar span has one component with two raw witnesses of weight $1/2$, so its apex is nonempty. The process $F_{\rm split}$ is also normalized, and its finite-measure trace is natively isomorphic to $X$. Writing the proposed extension as $\widetilde L$, an invertible feedback comparison would therefore require
\[
 LX\cong \widetilde L\bigl(\operatorname{Tr}^{\mathrm{fin}}_U(F_{\rm split})\bigr)
 \cong\operatorname{Tr}^{\mathrm{span}}_{LU}(LF_{\rm split}),
\]
The right-hand span has empty apex, contradicting the first assertion. Only this one split dilation is needed for the corollary; the packed dilation additionally proves dependence on presentation.
\end{proof}

\paragraph{Scope of the obstruction.}
This refutes arbitrary-dilation descent with the stated outputs and equality strength. It does not rule out a bridge retaining the normalized presentation, an additional chosen policy, a different comparison direction, or a changed span/effect interface. In particular a directed comparison from the empty traced span toward the nonempty image is not contradicted by the empty/nonempty argument; its existence and coherence must be proved separately. Nor does the result assert that every function or functor between the two semantic doctrines is impossible. It identifies the missing integration obligation: the atomic one-cell generation theorem does not make whole-law span admissibility an invariant of the closed finite-measure native process.

\end{document}
